\chapter{Definitions and elementary divisibility}
\label{ch:basic}

Everything in this project is built from the three definitions below. They are
project-local by design: Mathlib's \texttt{Nat.Prime}, the generic
\texttt{Irreducible}/\texttt{Prime} classes, and \texttt{Nat.gcd} are all out of
bounds, because reproducing Euclid's development from first principles is the
point of the project. Only elementary arithmetic and divisibility facts about
$\NN$ and $\ZZ$ are taken as given.

The informal source states each definition without the hypothesis $1 < p$. That
omission makes the definitions unusable, so we restore the hypothesis; the
discussion after Definition~\ref{def:is_irreducible} explains why.

\section{The three predicates}

\begin{definition}[prime]
  \leanok
  \label{def:is_prime}
  \lean{EuclidPrimes.IsPrime}
  % SOURCE: [euclid-notes], the first `def:` line (read from references/euclid-infinitude-of-primes.md)
  % SOURCE QUOTE: "def: A natural number p is prime if p | a * b => p | a or p | b"
  \textit{Source: euclid-infinitude-of-primes.md, first \texttt{def:} line.}
  A natural number $p$ is \emph{prime}, written $\Prm{p}$, when
  \[
    1 < p \quad\text{and}\quad \forall a, b \in \NN,\ p \mid a b \implies p \mid a \ \text{or}\ p \mid b .
  \]
\end{definition}

\begin{definition}[irreducible]
  \leanok
  \label{def:is_irreducible}
  \lean{EuclidPrimes.IsIrreducible}
  % SOURCE: [euclid-notes], the second `def:` line (read from references/euclid-infinitude-of-primes.md)
  % SOURCE QUOTE: "def: A natural number p is irreduicble if p = a * b => p = a or p = b"
  \textit{Source: euclid-infinitude-of-primes.md, second \texttt{def:} line.}
  A natural number $p$ is \emph{irreducible}, written $\Irr{p}$, when
  \[
    1 < p \quad\text{and}\quad \forall a, b \in \NN,\ p = a b \implies p = a \ \text{or}\ p = b .
  \]
\end{definition}

\paragraph{Why $1 < p$ is not optional.}
The source drops the unit condition from both definitions, and the resulting
predicates are false-friendly:
\begin{itemize}
  \item $p = 1$ satisfies the divisibility clause of
    Definition~\ref{def:is_prime} (since $1$ divides everything) and the
    factorisation clause of Definition~\ref{def:is_irreducible} (from
    $1 = ab$ over $\NN$ one gets $a = b = 1$), so $1$ would be both prime and
    irreducible.
  \item $p = 0$ satisfies the divisibility clause as well: $0 \mid ab$ means
    $ab = 0$, hence $a = 0$ or $b = 0$, i.e.\ $0 \mid a$ or $0 \mid b$. So $0$
    would be prime.
\end{itemize}
With either of those admitted, the factorisation theorem
(Theorem~\ref{thm:exists_factorization}) is false --- $[1,1,\dots]$ or a list
containing $0$ would be a legal ``factorisation'' of anything --- and the
infinitude argument collapses, because its final contradiction is exactly
$1 < q$ against $q \mid 1$. The condition $1 < p$ also supplies $p \neq 0$,
which every cancellation step in Chapter~\ref{ch:irreducible} needs.

\begin{definition}[greatest common divisor, by its universal property]
  \leanok
  \label{def:is_gcd}
  \lean{EuclidPrimes.IsGcd}
  % SOURCE: [euclid-notes], the third `def:` block (read from references/euclid-infinitude-of-primes.md)
  % SOURCE QUOTE: "def: if a b are naturals, the greatest common divisor gcd(a,b) = d is such that
  % - d | a
  % - d | b
  % - if c | a and c | b then c | d"
  \textit{Source: euclid-infinitude-of-primes.md, third \texttt{def:} block.}
  For $a, b, d \in \NN$ we say $d$ is \emph{a} greatest common divisor of $a$
  and $b$, written $\Gcd{a}{b}{d}$, when
  \[
    d \mid a, \qquad d \mid b, \qquad \forall c \in \NN,\ (c \mid a \ \text{and}\ c \mid b) \implies c \mid d .
  \]
\end{definition}

Note this is a \emph{relation}, not a function: the source writes
$\gcd(a,b) = d$, but nothing yet produces such a $d$, and nothing yet says it is
unique. Existence is Corollary~\ref{cor:exists_is_gcd} and uniqueness is
Lemma~\ref{lem:is_gcd_unique}, both in Chapter~\ref{ch:bezout}. Until then every
statement must quantify over ``some $d$ with $\Gcd{a}{b}{d}$''. Observe also
that the definition places no positivity constraint on $d$, and indeed
$\Gcd{0}{0}{0}$ holds; that degenerate case is harmless and is never excluded.

\section{A divisor is smaller than what it divides}

\begin{lemma}[strict divisor bound]
  \leanok
  \label{lem:dvd_lt_of_ne}
  \lean{EuclidPrimes.dvd_lt_of_ne}
  % SOURCE: [euclid-notes], the `lemma:` on distinct naturals greater than 1 (read from references/euclid-infinitude-of-primes.md)
  % SOURCE QUOTE: "lemma: let a,b be distinct natural numbers greater than 1. If a | b => a < b"
  \textit{Source: euclid-infinitude-of-primes.md, the third \texttt{lemma:}.}
  Let $a, b \in \NN$ with $1 < a$, $1 < b$ and $a \neq b$. If $a \mid b$ then
  $a < b$.
\end{lemma}
% SOURCE QUOTE PROOF: "proof: if a | b => exists a c with b = a * c. Since a b distinct, c > 1. Then
% b - a = a * c - a = a * (c - 1) > 0 => b - a > 0 => b > a"
\begin{proof}

\leanok
  The source argues via $b - a = a(c-1) > 0$. Over $\NN$ subtraction is
  truncated, so ``$b - a > 0$'' is not available as a step --- it is what we are
  trying to prove. We argue multiplicatively instead, which needs no subtraction
  at all.

  From $a \mid b$ choose $c \in \NN$ with $b = a c$. We rule out the two small
  values of $c$:
  \begin{itemize}
    \item If $c = 0$ then $b = 0$, contradicting $1 < b$.
    \item If $c = 1$ then $b = a$, contradicting $a \neq b$.
  \end{itemize}
  Hence $2 \le c$. Since $0 < a$, multiplying that inequality by $a$ gives
  $b = a c \ge 2a = a + a > a$, the last step because $0 < a$. So $a < b$.
\end{proof}

The next two lemmas package the arithmetic that the factorisation induction
(Chapter~\ref{ch:factorization}) needs: a non-trivial factor of $n$ is itself
bigger than $1$, and strictly smaller than $n$.

\begin{lemma}[a non-unit factor exceeds one]
  \leanok
  \label{lem:one_lt_factor}
  \lean{EuclidPrimes.one_lt_factor}
  Let $n, a, b \in \NN$ with $1 < n$ and $n = a b$. If $a \neq 1$ then $1 < a$.
\end{lemma}
\begin{proof}

\leanok
  If $a = 0$ then $n = 0 \cdot b = 0$, contradicting $1 < n$. So $a \neq 0$, and
  $a \neq 1$ by hypothesis; for a natural number this leaves exactly $1 < a$.
\end{proof}

\begin{lemma}[a proper factor is strictly smaller]
  \leanok
  \label{lem:factor_lt}
  \lean{EuclidPrimes.factor_lt}
  \uses{lem:dvd_lt_of_ne, lem:one_lt_factor}
  Let $n, a, b \in \NN$ with $1 < n$ and $n = a b$, and suppose $a \neq 1$ and
  $b \neq 1$. Then $a < n$.
\end{lemma}
\begin{proof}
  \uses{lem:dvd_lt_of_ne, lem:one_lt_factor}
  \leanok
  By Lemma~\ref{lem:one_lt_factor} applied to $n = a b$ we get $1 < a$, and
  applied to $n = b a$ we get $1 < b$.

  Next, $a \neq n$. Indeed, if $a = n$ then $n = n b$; since $1 < n$ gives
  $0 < n$, cancelling $n$ yields $b = 1$, contradicting $b \neq 1$.

  Finally $a \mid n$, because $n = a b$. Lemma~\ref{lem:dvd_lt_of_ne}, applied
  with the hypotheses $1 < a$, $1 < n$, $a \neq n$, gives $a < n$.
\end{proof}

By symmetry of the hypotheses ($n = b a$ with $b \neq 1$ and $a \neq 1$), the
same lemma also yields $b < n$; no separate statement is needed.

\section{Adequacy of the definitions}

The three predicates of the previous section are project-original in a strong
sense: the source's versions are defective --- they omit $1 < p$ --- and were
repaired here. Every later result in this project is a statement \emph{about}
$\Prm{\cdot}$, so a predicate that was accidentally too weak (satisfied by
numbers that are not prime) or too strong (failing on numbers that are) would
leave the whole development literally true and yet about the wrong notion. In
particular Theorem~\ref{thm:not_exists_prime_list} would still be a theorem.

The four results below close that gap. They pin $\Prm{\cdot}$ from three sides:
it excludes the units, it excludes a composite, it holds of $2$, and it agrees
with the divisor condition a reader has in mind. Each is proved \emph{from the
definitions alone} --- never from the development downstream, which would make
the check circular --- and none of them is used anywhere else, by design.

\begin{lemma}[units are not prime]
  \leanok
  \label{lem:not_is_prime_of_le_one}
  \lean{EuclidPrimes.not_isPrime_of_le_one}
  \uses{def:is_prime}
  Let $p \in \NN$ with $p \le 1$. Then $\Prm{p}$ fails. In particular neither
  $0$ nor $1$ is prime.
\end{lemma}
\begin{proof}
  \uses{def:is_prime}
  \leanok
  The first conjunct of Definition~\ref{def:is_prime} is $1 < p$, which
  contradicts $p \le 1$ outright.

  This is the formal counterpart of the discussion after
  Definition~\ref{def:is_irreducible}: both $0$ and $1$ satisfy the divisibility
  clause, so the exclusion is carried entirely by $1 < p$ and by nothing else.
\end{proof}

\begin{lemma}[two is prime]
  \leanok
  \label{lem:is_prime_two}
  \lean{EuclidPrimes.isPrime_two}
  \uses{def:is_prime}
  $\Prm{2}$ holds.
\end{lemma}
\begin{proof}
  \uses{def:is_prime}
  \leanok
  The first conjunct, $1 < 2$, is immediate. For the second, let $a, b \in \NN$
  with $2 \mid ab$, and suppose $2 \nmid a$; we must produce $2 \mid b$.

  Divisibility by $2$ is the vanishing of the residue mod $2$, so the hypotheses
  read $ab \bmod 2 = 0$ and $a \bmod 2 \neq 0$; as $a \bmod 2 < 2$, the latter
  forces $a \bmod 2 = 1$. Residues are multiplicative,
  \[
    ab \bmod 2 = \bigl( (a \bmod 2)(b \bmod 2) \bigr) \bmod 2
      = (b \bmod 2) \bmod 2 = b \bmod 2 ,
  \]
  the last step because $b \bmod 2 < 2$. Hence $b \bmod 2 = ab \bmod 2 = 0$,
  i.e.\ $2 \mid b$.

  Note this is a parity argument and uses nothing but multiplicativity of
  residues: it does not go through irreducibility, so it does not depend on
  Chapter~\ref{ch:bezout} or Chapter~\ref{ch:irreducible}. That independence is
  the point --- a witness that the definition is satisfiable must not rest on the
  development it is meant to validate.
\end{proof}

\begin{lemma}[four is not prime]
  \leanok
  \label{lem:not_is_prime_four}
  \lean{EuclidPrimes.not_isPrime_four}
  \uses{def:is_prime}
  $\Prm{4}$ fails.
\end{lemma}
\begin{proof}
  \uses{def:is_prime}
  \leanok
  Suppose $\Prm{4}$. Since $4 = 2 \cdot 2$ we have $4 \mid 2 \cdot 2$, so the
  divisibility clause of Definition~\ref{def:is_prime} applied with
  $a = b = 2$ yields $4 \mid 2$ (either disjunct says exactly that). A positive
  number is at least as large as any of its divisors, so $4 \le 2$ --- absurd.

  Together with Lemma~\ref{lem:not_is_prime_of_le_one} and
  Lemma~\ref{lem:is_prime_two} this rules out the two ways the definition could
  be wrong by inclusion, and exhibits one number it does admit.
\end{proof}

\begin{lemma}[irreducibility is the condition on divisors]
  \leanok
  \label{lem:is_irreducible_iff_forall_dvd}
  \lean{EuclidPrimes.isIrreducible_iff_forall_dvd}
  \uses{def:is_irreducible}
  For $p \in \NN$,
  \[
    \Irr{p} \iff \Bigl(1 < p \ \text{ and }\ \forall d \in \NN,\
      d \mid p \implies (d = 1 \ \text{ or }\ d = p)\Bigr).
  \]
\end{lemma}
\begin{proof}
  \uses{def:is_irreducible}
  \leanok
  The $1 < p$ conjunct is common to both sides, so only the two quantified
  clauses need relating.

  ($\Rightarrow$) Assume $\Irr{p}$ and let $d \mid p$, say $p = dc$. Applying
  the factorisation clause to $p = dc$ gives $p = d$ or $p = c$. In the first
  case $d = p$ and we are done. In the second, substituting $c = p$ into
  $p = dc$ gives $p = dp$, that is $p \cdot 1 = p \cdot d$; since $1 < p$ gives
  $p \neq 0$, cancelling $p$ yields $d = 1$.

  ($\Leftarrow$) Assume the divisor condition and let $p = ab$. Then $a \mid p$,
  so $a = 1$ or $a = p$. If $a = p$ then $p = a$, the left disjunct. If $a = 1$
  then $p = 1 \cdot b = b$, the right disjunct. (This direction needs no
  positivity: it never cancels.)
\end{proof}

\paragraph{What this buys.} The right-hand side of
Lemma~\ref{lem:is_irreducible_iff_forall_dvd} is the schoolbook definition of a
prime number: bigger than $1$, and divisible by nothing except $1$ and itself.
Composing it with Corollary~\ref{cor:prime_iff_irreducible} --- which is where
the mathematical content of the project's first half sits --- transports that
characterisation from $\Irr{\cdot}$ to $\Prm{\cdot}$, so the predicate that
Corollary~\ref{cor:infinite_setOf_prime} proves to have infinitely many
witnesses is provably the familiar one, and not merely a predicate that
resembles it. That composition is
Corollary~\ref{cor:prime_iff_forall_dvd}; it is a stated and proved result of
Chapter~\ref{ch:irreducible} rather than a remark here, since
Corollary~\ref{cor:prime_iff_irreducible} is not available this early.

Note the direction of the dependency. Nothing on the path to Euclid's theorem
uses a lemma of this section: its only consumer anywhere is
Corollary~\ref{cor:prime_iff_forall_dvd}, and that corollary in turn is used by
nothing --- every result of
Chapters~\ref{ch:bezout}--\ref{ch:infinitude} is proved without it. The adequacy
results are a check on the definitions, never a step in the argument.
